\documentclass[10pt,a4paper]{amsart}
\usepackage[margin=19mm]{geometry}
\usepackage{amsmath,amssymb,mathtools}
\usepackage[scr=rsfs,cal=euler]{mathalfa}
\usepackage{xfrac}
\usepackage[unicode,hidelinks,bookmarks=false]{hyperref}
\newcommand{\ie}{i.e.\ }
\newcommand{\cf}{cf.\ }
\newcommand{\resp}{respectively\ }
\newcommand{\Subsection}{\subsection}
\providecommand{\nobreakdash}{\nobreak\hbox{-}}
\setlength{\emergencystretch}{2em}
\multlinegap0pt
\usepackage{bm}
\usepackage{bbm}
\usepackage{dsfont}
\usepackage{xspace}
\usepackage{cancel}
\usepackage{mathtools}
\usepackage{amsthm}
\makeatletter
\@ifundefined{theorem}{\newtheorem{theorem}{Theorem}}{}
\@ifundefined{proposition}{\newtheorem{proposition}[theorem]{Proposition}}{}
\makeatother
\newcommand{\im}{\mathrm{im}}
\title[Smooth relative connections on quiver bundles]{Smooth relative connections on quiver bundles}
\author{Pavan Adroja, Sanjay Amrutiya, Riddhi Patil}
\date{Source: arXiv:2602.18781v1 (CC BY 4.0)}
\begin{document}
\maketitle

\noindent\emph{Source and licence.} Smooth relative connections on quiver bundles. Version arXiv:2602.18781v1 (CC BY 4.0), distributed under the Creative Commons Attribution 4.0 International licence, as recorded in the arXiv licence metadata for this exact version and shown on the abstract page https://arxiv.org/abs/2602.18781v1. Full rights notice: licenses/arxiv-2602.18781-CC-BY-4.0.txt.\par\smallskip

\noindent\textit{Editorial scope.} Counted unit: the Proposition whose source label is \texttt{tree} in the article (source0.tex lines 641--643, source label \texttt{tree}), delivered together with the article's own complete original proof of it (source0.tex lines 644--907, one \texttt{proof} environment of 264 lines). No mathematics is written, repaired or re-derived: statement and proof are the article's own text line for line. Required context retained uncounted: none. Every definition, display and label of those lines is retained unchanged. Omitted from this package and not redistributed: 1--640, 908--1071. Nothing inside a retained excerpt is omitted or altered: every retained body is delivered exactly as the article's lines are written, so no omission receipt of any kind accompanies this package. Citations: the retained excerpts carry no citation command at all, so no bibliography entry of the article is transcribed and no reference is redistributed. Reference adaptation: none was needed and none was made; every reference inside the delivered unit resolves inside it, so no unresolved reference or citation is produced. Printed slips kept literal: no printed word, symbol, hypothesis or number of the counted statement or of its proof was altered. Layout and class: the article's own class is \texttt{amsart} with packages including amsfonts, amssymb, hyperref, graphicx, epsfig, latexsym, amsmath, amsthm, calligra, mathrsfs, xy, comment, xcolor, tikz-cd; the wrapper is the lane's pinned \texttt{amsart} editorial wrapper at 10pt on A4, which restates the theorem-like environments the retained text needs and adds the article's own macro definitions that the retained text needs (\texttt{\textbackslash im}), with extra wrapper packages (bm, bbm, dsfont, xspace, cancel, mathtools). The delivered numbering is the unit's own. Assets: no retained span contains a figure environment, a graphics-inclusion command, a PDF-page inclusion, a table or a TikZ picture; no image file is copied into this package, the package declares no asset dependency and no image file is redistributed with it. Licence: version arXiv:2602.18781v1 of this article is distributed under the Creative Commons Attribution 4.0 International licence, as recorded in the arXiv licence metadata for this exact version and shown on the abstract page https://arxiv.org/abs/2602.18781v1. A full rights notice is delivered at licenses/arxiv-2602.18781-CC-BY-4.0.txt. Characters: every retained excerpt is pure ASCII, so the delivered engine, XeLaTeX with its default Latin Modern text font, reports no missing character.
\begin{proposition}\label{tree}
Let $Q$ be a rooted directed tree with root $r$. A $Q$-bundle $\mathcal{R} = (\mathcal{E}, \phi)$ admits a smooth relative connection if and only if every path morphism $\phi_p$ in $Q$ is of constant rank.
\end{proposition}

\begin{proof}
Suppose that, for every path $p$ in $Q$, the associated morphism $\phi_p$ is of constant rank. Let $v_1, \dots, v_l$ be the leaves of $Q$. 

For each leaf $v_i$, let $p_{v_i} = a_{n_i i} \cdots a_{1 i}$
be the unique path from $r$ to $v_i$. 
For $1 \le k \le n_i$, denote by
$p_{v_i}^{(k)} := a_{k i} \cdots  a_{1 i}$
the initial subpath of length $k$.
As in the $\mathbb{A}_n$ case, the constant rank condition gives a filtration by smooth subbundles
\[
\ker(\phi_{p_{v_i}^{(1)}})
\subseteq
\ker(\phi_{p_{v_i}^{(2)}})
\subseteq
\cdots
\subseteq
\ker(\phi_{p_{v_i}^{(n_i)}}) \subseteq
\mathcal{E}_r\,.
\]

Choose smooth complementary subbundles representing the successive quotients. 
More precisely, choose subbundles
\[
\mathcal{E}_r^{(v_i,0)}, \dots, \mathcal{E}_r^{(v_i,n_i)} \subset \mathcal{E}_r
\]
such that
\[
\ker(\phi_{p_{v_i}^{(1)}}) = \mathcal{E}_r^{(v_i,0)},
\]
\[
\ker(\phi_{p_{v_i}^{(k+1)}})
=
\ker(\phi_{p_{v_i}^{(k)}})
\oplus
\mathcal{E}_r^{(v_i,k)}
\quad \text{for } 1 \le k < n_i,
\]
and
\[
\mathcal{E}_r
=
\ker(\phi_{p_{v_i}^{(n_i)}})
\oplus
\mathcal{E}_r^{(v_i,n_i)}.
\]
This yields a direct sum decomposition
\[
\mathcal{E}_r
=
\bigoplus_{m_i=0}^{n_i}
\mathcal{E}_r^{(v_i,m_i)}.
\]
We then define the final decomposition of $\mathcal{E}_r$ by
\[
\mathcal{E}_r
=
\bigoplus_{(m_1,\dots,m_l)}
\left(
\mathcal{E}_r^{(v_1,m_1)}
\cap
\cdots
\cap
\mathcal{E}_r^{(v_l,m_l)}
\right),
\]
where $0 \le m_i \le n_i$ for each $i$.

Now fix a vertex $u \neq r$. Since $Q$ is a tree, there is a unique path $p_u$ from $r$ to $u$. Write
$p_u = a_q \cdots a_1.$
Then $\im(\phi_{p_u}) \subseteq \mathcal{E}_u$.
For each $j = 1, \dots, q$, define the intermediate subpaths ending at $u$ by
\[
p_u^{(j)} := a_q \cdots a_j,
\]
and set
\[
I_u^{(j)} := \im(\phi_{a_q} \circ \cdots \circ \phi_{a_j})
\subseteq \mathcal{E}_u.
\]
These are precisely the images of the path morphisms corresponding to subpaths of $p_u$ that end at $u$. Since all path morphisms have constant rank, each $I_u^{(j)}$ is a smooth subbundle of $\mathcal{E}_u$, and the inclusions
\[
I_u^{(1)} \subseteq \cdots \subseteq I_u^{(q)} \subseteq \mathcal{E}_u
\]
form a filtration by subbundles. Choose smooth complementary subbundles
\[
J_u^{(0)} \subseteq I_u^{(1)}, \qquad
J_u^{(j)} \subseteq I_u^{(j+1)} \ \text{for } 1 \le j < q, \qquad \text{and} \qquad J_u^{(q)} \subseteq \mathcal{E}_u,
\]
such that
\[
I_u^{(1)} = J_u^{(0)}, \qquad
I_u^{(j+1)} = I_u^{(j)} \oplus J_u^{(j)} \ \text{for } 1 \le j < q, \qquad \text{and} \qquad
\mathcal{E}_u = I_u^{(q)} \oplus J_u^{(q)}.
\]
This yields a decomposition
\[
\mathcal{E}_u
=
\bigoplus_{j=0}^{q}
J_u^{(j)}.
\]


Now let $w_1, \dots, w_s$ be the leaves of the rooted directed subtree with root $u$. We proceed as in the root case. For each leaf $w_i$, let
$p_{w_i} = a_{e_i i} \cdots a_{1 i}$
be the unique path from $u$ to $w_i$. 
For $1 \le k \le e_i$, denote by
\[
p_{w_i}^{(k)} := a_{k i} \cdots a_{1 i}
\]
the initial subpath of length $k$.

Since all path morphisms have constant rank, the kernels
\[
\ker(\phi_{p_{w_i}^{(1)}})
\subseteq
\cdots
\subseteq
\ker(\phi_{p_{w_i}^{(e_i)}})
\]
form a filtration by subbundles of $\mathcal{E}_u$. Choose smooth splittings of this filtration and let 
$\mathcal{E}_u^{(w_i,k)}$ be complementary subbundles representing the successive quotients.
\[
\mathcal{E}_u
=
\bigoplus_{k=0}^{e_i}
\mathcal{E}_u^{(w_i,k)}.
\]
Taking intersections over all leaves of the subtree, we obtain
\[
\mathcal{E}_u
=
\bigoplus_{(k_1,\dots,k_s)}
\left(
\mathcal{E}_u^{(w_1,k_1)}
\cap
\cdots
\cap
\mathcal{E}_u^{(w_s,k_s)}
\right).
\]

Finally, refining this decomposition with the image filtration
\[
I_u^{(1)} \subseteq \cdots \subseteq I_u^{(q)} \subseteq \mathcal{E}_u,
\]
we obtain the final decomposition
\[
\mathcal{E}_u
=
\bigoplus_{j=0}^{q}
\;
\bigoplus_{(k_1,\dots,k_s)}
\left(
\mathcal{E}_u^{(w_1,k_1)}
\cap
\cdots
\cap
\mathcal{E}_u^{(w_s,k_s)}
\cap
J_u^{(j)}
\right).
\]


We now construct a smooth relative connection on the $Q$-bundle.

At the root $r$, consider the refined decomposition
\[
\mathcal{E}_r
=
\bigoplus_{(m_1,\dots,m_l)}
\left(
\mathcal{E}_r^{(v_1,m_1)}
\cap
\cdots
\cap
\mathcal{E}_r^{(v_l,m_l)}
\right).
\]
Choose an arbitrary smooth connection on each summand and denote the
resulting connection on $\mathcal{E}_r$ by $\nabla_{\mathcal{E}_r}$.

Since $Q$ is a rooted directed tree with root $r$, every vertex $v \neq r$ has a unique
incoming arrow. We now define the connection on each vertex by
moving outward from the root along the arrows.

Let $a : u \to v$ be an arrow such that a connection has already
been defined on $\mathcal{E}_u$. Write the refined decomposition of $\mathcal{E}_u$ as
\[
\mathcal{E}_u
=
\bigoplus_{j=0}^{q}
\;
\bigoplus_{(k_1,\dots,k_s)}
\left(
\mathcal{E}_u^{(w_1,k_1)}
\cap
\cdots
\cap
\mathcal{E}_u^{(w_s,k_s)}
\cap
J_u^{(j)}
\right),
\]
where the summands are those obtained from the intersection of
the kernel and image filtrations.

For each such summand
\[
\mathcal{S}_u
=
\left(
\mathcal{E}_u^{(w_1,k_1)}
\cap
\cdots
\cap
\mathcal{E}_u^{(w_s,k_s)}
\cap
J_u^{(j)}
\right),
\]
the restriction
\[
\phi_a\big|_{\mathcal{S}_u}
:
\mathcal{S}_u
\longrightarrow
\mathcal{E}_v
\]
is either zero or an isomorphism onto its image, and this image
is a direct sum of certain summands in the corresponding
decomposition of $\mathcal{E}_v$.

If $\phi_a|_{\mathcal{S}_u} = 0$, we choose an arbitrary smooth
connection on the corresponding summands of $\mathcal{E}_v$. If $\phi_a|_{\mathcal{S}_u}
:
\mathcal{S}_u
\longrightarrow
\phi_a(\mathcal{S}_u)$ is an isomorphism, we define the connection on
$\phi_a(\mathcal{S}_u)$ by :
\[
\nabla_{\phi_a(\mathcal{S}_u)}
=
(\mathrm{id}_{T_X^*}\otimes \phi_a) \circ \nabla_{\mathcal{S}_u} \circ \phi_a^{-1}.
\]
Taking the direct sum over all summands defines a smooth
connection on $\mathcal{E}_v$.

Proceeding along all arrows starting from the root,
this defines smooth connections on every vertex of $Q$.
Since each vertex has a unique incoming arrow,
the construction is well defined.

By construction, for every arrow $a : u \to v$ we have
\[
\nabla_{\mathcal{E}_v} \circ \phi_a
=
(\mathrm{id}_{T^*_X} \otimes \phi_a)
\circ
\nabla_{\mathcal{E}_u}.
\]
Hence, the $Q$-bundle admits a smooth relative connection.
\end{proof}

\end{document}
